\documentclass[10pt]{article}
\usepackage[utf8]{inputenc}
\usepackage[T1]{fontenc}
\usepackage[margin=0.78in]{geometry}
\usepackage{amsmath,amssymb}
\usepackage{booktabs}
\usepackage{xcolor}
\usepackage{enumitem}
\usepackage[small,skip=3pt]{caption}
\usepackage[colorlinks=true,linkcolor=blue,citecolor=blue]{hyperref}
\setlength{\parskip}{2pt}
\newcommand{\code}[1]{\texttt{\small #1}}
\pagestyle{empty}

\title{\vspace{-2.2em}\textbf{Reproducing, Diagnosing, and Rescuing the CovidGAN Claim}\\[2pt]
\large --- Two-Page Summary ---}
\author{\normalsize Hila Fishman \and \normalsize Maya Hayat}
\date{}

\begin{document}
\maketitle
\vspace{-3.2em}

\section*{The claim under test}
Waheed et al.\ (2020) propose \textbf{CovidGAN}: an AC-GAN synthesizes chest X-rays, and those synthetic
images augment a VGG16 COVID-19 detector. Their headline claim is a \textbf{+10-point accuracy lift}
(85\%~$\to$~95\%) on a 192-image real test set, concentrated in COVID recall. We reconstructed the pipeline in
PyTorch, tried to reproduce the claim, diagnosed why it failed, fixed it, and re-ran everything on a second
independent dataset to check our conclusions were not artifacts of our data.

\section*{Stage 1 --- the claim does not reproduce}
Our faithful reconstruction (2000-epoch GAN, frozen VGG16 head, paper's split and hyperparameters)
produced \textbf{two anomalies}: the baseline was \emph{higher} than the paper's (90.62\% vs.\ 85\%), and
augmentation was \textbf{flat} ($-0.52$), not $+10$. Five controlled tests explain both:

\begin{enumerate}[nosep,leftmargin=1.4em]
  \item \textbf{Not a bug.} A parameter-count and layer-freeze audit matched the paper exactly.
  \item \textbf{Not a cross-source shortcut.} The paper mixes IEEE (COVID) and Kaggle (Normal) sources, so
        we rebuilt the split with both classes from one source. Accuracy was unchanged ($\approx$91\%).
  \item \textbf{Not a coarse shortcut.} Blurring real images to $4{\times}4$ collapses accuracy by 16 points
        and halves COVID recall --- the model reads genuine fine detail, not global brightness or framing.
  \item \textbf{The fakes carry almost no pathology.} Trained on synthetic images alone, the classifier
        scores \textbf{65.45\%} on real data --- barely above the 62.5\% majority-class floor. The AC-GAN
        stamps a class-conditioned \emph{brightness fingerprint}, plainly visible in the samples, rather
        than learning localized pathology.
  \item \textbf{Quality is not the blocker.} FID halved over training (504~$\to$~273) with \emph{zero}
        downstream change.
\end{enumerate}

\noindent\textbf{Diagnosis.} The high baseline reflects \emph{data drift} --- today's Kaggle database is
cleaner and more separable than the paper's March-2020 snapshot, leaving little headroom. The flat
augmentation reflects the paper's \emph{frozen backbone}: with only 33K trainable parameters, synthetic
data can nudge a linear boundary but cannot reshape features.

\section*{Stage 2 --- two tracks}
\textbf{Track A (generator).} The paper samples $z\sim\mathcal{N}(0,0.02)$, which leaves the noise nearly
constant so the class label dominates --- the direct cause of the fingerprint. We restored
$z\sim\mathcal{N}(0,1)$ and added DiffAugment, spectral normalization, and a projection discriminator
(removing the auxiliary head that rewards a cheap class marker).
\textbf{Track B (classifier).} We unfroze the top VGG16 blocks with a BatchNorm head and discriminative
learning rates (head $10^{-3}$, backbone $10^{-5}$), giving augmentation something to act on.

\begin{table}[htbp]\centering
\caption{\textbf{Paper vs.\ reconstruction vs.\ improved.} Same 192-image real test set. Stage 2 rows are
mean $\pm$ std over 5 seeds.}
\begin{tabular}{@{}lcccc@{}}
\toprule
\textbf{Model} & \textbf{Trainable} & \textbf{CNN-AD} & \textbf{CNN-SA} & \textbf{Lift} \\
\midrule
Paper (Waheed et al.) & $\sim$33K & 85.00\% & 95.00\% & $+10.00$ \\
Stage 1 reconstruction (frozen) & $\sim$33K & 90.62\% & 90.10\% & $-0.52$ \\
Stage 2 improved ($uf{=}1$) & 7.11M & 92.50\% $\pm$ 0.71 & 93.65\% $\pm$ 1.16 & $+1.15$ \\
\textbf{Stage 2 improved ($uf{=}2$)} & 13.01M & \textbf{94.48\% $\pm$ 0.97} & \textbf{96.67\% $\pm$ 0.71} & $\mathbf{+2.19}$ \\
\bottomrule
\end{tabular}
\end{table}

\noindent Unfreezing raised the baseline \emph{and} restored a real, multi-seed augmentation lift, with
CNN-SA passing the paper's 95\%. The paper's claim reproduces \textbf{in kind, not in magnitude}.

\section*{The central finding --- generator quality is inversely related to downstream value}
Feeding four different synthetic pools to the \emph{same} classifier, changing nothing else, gives a result
that runs opposite to expectation: as the generator improves by every measure available for judging it
alone, the downstream benefit \textbf{decays to nothing}.

\begin{table}[htbp]\centering
\caption{Generator-side quality (left) versus downstream value (right, $uf{=}2$, 5 seeds, against a common
CNN-AD of 94.48\% $\pm$ 0.97). Transfer probe = classifier trained on synthetic alone, tested on real;
floor $=62.5\%$.}
\begin{tabular}{@{}lccccc@{}}
\toprule
\textbf{Generator} & \textbf{Ep.} & \textbf{FID $\downarrow$} & \textbf{Transfer $\uparrow$} & \textbf{CNN-SA} & \textbf{Lift} \\
\midrule
Stage-1 AC-GAN (baseline) & 2000 & 272.7 & 65.45\% & 96.67\% & $+2.19$ \\
Improved AC-GAN & 300 & 225.7 & 72.40\% & \textbf{96.98\%} & $\mathbf{+2.50}$ \\
Projection discriminator & 300 & 155.0 & 74.65\% & 96.04\% & $+1.56$ \\
Projection discriminator & 600 & \textbf{121.3} & \textbf{82.29\%} & 94.79\% & $+0.31$ \emph{(n.s.)} \\
\bottomrule
\end{tabular}
\end{table}

\noindent The best generator we built produces the \emph{only} lift indistinguishable from zero. A
matched-epoch control rules out training budget as the explanation: retrained at 300 epochs, the baseline
architecture reaches $+1.77$, so neither improved generator significantly separates from it
($+0.73$, $t\approx1.20$; $-0.21$, $t\approx0.24$).

\section*{Capacity, not generator quality, is the governing variable}
Crossing every pool against both classifier regimes (3 seeds frozen, 5 unfrozen): under the \textbf{frozen}
head \emph{no} pool achieves a positive lift --- the best manages $+0.17$, and two significantly
\emph{harm} the classifier ($-1.39$ each). Under the \textbf{unfrozen} encoder the identical pools all help
($+0.31$ to $+2.50$). Moving frozen~$\to$~unfrozen is worth 5--6 points on every row; moving between
generators is worth at most 2. Under scarcity this sharpens: with the frozen head and our \emph{best} pool,
augmentation loses $-1.91$ points at full data and $-8.16$ at 10\% real, where 93 real images face 3{,}068
synthetic ones.

\section*{Second dataset --- validation on CBIS-DDSM}
Because our leading explanation for the high baseline was COVID-specific data drift, we re-ran the
pipeline unchanged on \textbf{CBIS-DDSM} mammography, a static, versioned benchmark the drift argument
cannot reach. Every checked conclusion held: the naive setup failed there too
(64.02\%~$\to$~65.34\%, flat); unfreezing raised the baseline (66.00~$\to$~69.45\%) and produced its first
reproducible augmentation win (70.77\%, both seeds); and a better generator again degraded the
high-capacity classifier (66.67\%). \textbf{The failure to reproduce is therefore a property of the method
under these conditions, not an artifact of our COVID data.} The two datasets diverge on one point: at
matched epochs DDSM's improved GAN gains $+6.9$ points where ours gains nothing --- consistent with
headroom, since DDSM sits near 64\% and our classifier near 94.5\%.

\section*{What we learned}
Two knobs govern whether synthetic augmentation helps: \textbf{capacity} (a frozen extractor stays flat or
gets worse, whatever it is fed) and \textbf{headroom} (augmentation pays when data-starved, little near
ceiling). The methodological lesson is sharpest: \textbf{FID and transfer probes were inversely related to
downstream value across our four generators.} Judging a generator alone is an unsafe proxy --- only a direct
AD-vs-SA comparison tests an augmentation claim. The paper's finding is real but \emph{contingent}: it needs
a classifier with capacity to respond and a baseline not already at ceiling.

\vspace{2pt}
\noindent\footnotesize\textbf{Reproducibility.} All runs are seeded and reproduce exactly; every number is
written to a \code{summary.json} recording seed, unfreeze depth, and pool, so each table cell traces to the
run that produced it. The 2000-epoch GAN took $\approx$18h on an Apple-silicon MPS GPU; each detector
trains in about a minute, which made the multi-seed grids practical.

\end{document}
